As AI agents act with less supervision, the people they affect need to be able to state their needs, correct the agent and see what it did. We present the Multi-stakeholder Relational Agency Framework (MSRAF), a gym that scores what an agent does when several affected people want different things. MSRAF reports three metrics separately and never sums them: Cooperation (hearing and including every party's needs), Corrigibility (leaving real room to say no and changing course without defending itself, while still using judgment) and Candor (making reasons, limits and consequential problems visible). Each targets an individual-action version of failures that arguments about catastrophic AI risk describe: resisting correction, deception, and acting without authority or consent. One LLM judge scores 18 checks, two of them computed from logs, from a rubric not yet human-reviewed. In this revised pilot, four models acted as facilitator in 48 episodes on one handcrafted three-stakeholder scenario, in a two-by-two prompt experiment (a Convergent Facilitation toolkit and an active participant framing). The toolkit raised Cooperation by 0.18 (95 percent CI 0.10 to 0.26), partly through hidden-need checks that its text also instructs, and in two of three scenarios overall. Corrigibility was at ceiling, so its checks need correction-pressure events. Candor mostly tracked whether agents finished a report before the turn cap, and a second judge disagreed often on it. The framing had no detectable effect. Clean results in toy scenarios would not establish safety.
